\chapter{Véc-tơ và hệ trục tọa độ trong không gian}

\section{Véc-tơ trong không gian}

\exercise Cho hình hộp $ABCD.A'B'C'D'$.
\begin{enumerate}
   \item Biểu diễn véc-tơ $\overrightarrow{AC'}$ theo ba véc-tơ $\overrightarrow{AB}, \overrightarrow{AD}, \overrightarrow{AA'}$.
   \item Rút gọn $S = \overrightarrow{AB} + \overrightarrow{B'C'} + \overrightarrow{DD'}$.
\end{enumerate}

\exercise Cho hình lập phương $ABCD.A'B'C'D'$ có cạnh bằng $a$.
\begin{enumerate}
   \item Tính tích vô hướng $\overrightarrow{AB} \cdot \overrightarrow{A'C'}$.
   \item Tính góc giữa hai véc-tơ $\overrightarrow{AB'}$ và $\overrightarrow{A'C'}$.
\end{enumerate}

\exercise Cho tứ diện $ABCD$. Gọi $G$ là trọng tâm tam giác $BCD$.
\begin{enumerate}
   \item Biểu diễn véc-tơ $\overrightarrow{AG}$ theo ba véc-tơ $\overrightarrow{AB}, \overrightarrow{AC}, \overrightarrow{AD}$.
   \item Gọi $M$ là trung điểm cạnh $AB$, $N$ là trung điểm cạnh $CD$ và $I$ là trung điểm của đoạn $MN$. Chứng minh ba điểm $A, I, G$ thẳng hàng.
\end{enumerate}

\section{Hệ trục tọa độ trong không gian}

\exercise Trong không gian $Oxyz$, cho hai điểm $A(1; -2; 3)$ và $B(3; 0; -1)$.
\begin{enumerate}
   \item Tìm tọa độ của véc-tơ $\overrightarrow{AB}$ và tính độ dài đoạn thẳng $AB$.
   \item Tìm tọa độ hình chiếu vuông góc $A'$ của điểm $A$ trên mặt phẳng $(Oxy)$ và hình chiếu $B'$ của $B$ trên trục $Oz$.
\end{enumerate}

\exercise Trong không gian $Oxyz$, cho ba điểm $A(2; 1; -1)$, $B(0; 3; 1)$, $C(1; -1; 2)$.
\begin{enumerate}
   \item Tìm tọa độ điểm $D$ sao cho tứ giác $ABCD$ là hình bình hành.
   \item Tìm tọa độ điểm $M$ thuộc trục $Oy$ sao cho $M$ cách đều hai điểm $A$ và $B$.
\end{enumerate}

\exercise Trong không gian $Oxyz$, cho tam giác $ABC$ có $A(1; 2; -1)$, $B(2; 3; 1)$ và $C(0; 1; 2)$.
\begin{enumerate}
   \item Tìm tọa độ trọng tâm $G$ của tam giác $ABC$ và tọa độ điểm $E$ sao cho $C$ là trung điểm của $AE$.
   \item Gọi $A_1, B_1, C_1$ lần lượt là hình chiếu vuông góc của $A, B, C$ lên mặt phẳng $(Oyz)$. Tính diện tích tam giác $A_1B_1C_1$.
\end{enumerate}

\section{Biểu thức toạ độ của các phép toán véc-tơ}

\exercise Trong không gian $Oxyz$, cho hai véc-tơ $\vec{u} = (2; -1; 1)$ và $\vec{v} = (1; 3; -2)$.
\begin{enumerate}
   \item Tìm tọa độ của véc-tơ $\vec{w} = 2\vec{u} - 3\vec{v}$.
   \item Tính tích vô hướng $\vec{u} \cdot \vec{v}$.
\end{enumerate}

\exercise Trong không gian $Oxyz$, cho ba điểm $A(1; 0; 1)$, $B(2; 1; 3)$ và $C(3; 2; 1)$.
\begin{enumerate}
    \item Tính cos của góc $\widehat{BAC}$ trong tam giác $ABC$.
    \item Tìm tọa độ véc-tơ $\vec{a}$ cùng hướng với véc-tơ $\overrightarrow{AB}$ và có độ dài bằng $3\sqrt{6}$.
\end{enumerate}

\exercise Trong không gian $Oxyz$, cho hai điểm $A(2; -1; 3)$ và $B(0; 3; 1)$.
\begin{enumerate}
   \item Tìm tọa độ điểm $C$ thuộc mặt phẳng $(Oxy)$ sao cho tam giác $ABC$ cân tại $C$ và hoành độ của $C$ bằng $3$.
   \item Tìm điểm $M$ thuộc trục $Ox$ sao cho $\overrightarrow{MA} \cdot \overrightarrow{MB} = 3$.
\end{enumerate}

\exercise Trong một bài tập mô phỏng bay không gian, vị trí của hai máy bay mô hình tại thời điểm $t$ phút ($t \ge 0$) được xác định bởi tọa độ:
$$M_1(1 + 2t; 2 - t; 3 + t) \quad \text{và} \quad M_2(4 + t; 1 + t; 1 + 3t)$$
\begin{enumerate}
    \item Tính khoảng cách giữa hai máy bay tại thời điểm $t = 2$ phút.
    \item Tìm thời điểm $t$ để khoảng cách giữa hai máy bay là nhỏ nhất và tính khoảng cách nhỏ nhất đó.
\end{enumerate}